\documentclass[11pt,a4paper]{article}
\usepackage{brailleexamen}
\usepackage[defaultfam]{montserrat}
\usepackage[spanish,es-noshorthands,es-tabla]{babel}
\hypersetup{hidelinks}
\setlength{\parindent}{0pt}
\setlength{\emergencystretch}{1em}
\begin{document}
\begin{tabularx}{\textwidth}{@{}Xr@{}}
{\color{primary}\Large\bfseries IES LUIS BRAILLE} & {\color{secondary}\bfseries BASES DE DATOS · RA3} \\
Departamento de Informática · Coslada & Curso 2026/2027 \\
\end{tabularx}
{\color{primary}\rule{\textwidth}{1.2pt}}
\begin{center}{\Large\bfseries Simulacro de consultas SELECT}\end{center}
\vspace{0.2cm}
\begin{tabularx}{\textwidth}{|p{2.4cm}|X|p{2.2cm}|X|}
\hline
\textbf{Módulo:} & Bases de Datos & \textbf{Curso:} & 1.º DAM/DAW \\ \hline
\textbf{Nombre:} & & \textbf{Apellidos:} & \\ \hline
\textbf{Fecha:} & & \textbf{Duración:} & 75 minutos \\ \hline
\textbf{Nota:} & \multicolumn{3}{l|}{\hspace{2cm}/ 10 puntos} \\ \hline
\end{tabularx}
\vspace{0.35cm}
\begin{tcolorbox}[title=Instrucciones]
Resuelve las ocho consultas sobre la base de datos \texttt{tienda\_online}. Entrega un único archivo \texttt{.sql} con las consultas numeradas. Cada consulta es independiente. Ordena exactamente como se pide. \textbf{Consultas 1 a 6: 1 punto cada una; consultas 7 y 8: 2 puntos cada una.} Las funciones indicadas en las preguntas 7 y 8 son obligatorias.
\end{tcolorbox}
{\small\textbf{Tablas disponibles:} \texttt{clientes}, \texttt{productos}, \texttt{pedidos}, \texttt{detalle\_pedido} y \texttt{pagos}. Puedes consultar su estructura con \texttt{DESCRIBE}.}
\par\vspace{0.25cm}
\textbf{Consulta 1 (1 punto).} Muestra las diez líneas de pedido con mayor importe (cantidad por precio\_unitario). Incluye identificador de línea, cantidad, precio unitario e importe; desempata por identificador.\par\vspace{0.65em}
\textbf{Consulta 2 (1 punto).} Muestra el identificador, fecha y año de pago de los quince primeros pagos por identificador. Conserva los pagos sin fecha.\par\vspace{0.65em}
\textbf{Consulta 3 (1 punto).} Construye una etiqueta para cada cliente con el formato «nombre \textless{}email\textgreater{}». Muestra el identificador y la etiqueta, ordenados por la quinta letra del nombre.\par\vspace{0.65em}
\textbf{Consulta 4 (1 punto).} Suma el coste de los pedidos por número de mes, reuniendo todos los años. Muestra mes e importe, ordenados por mes.\par\vspace{0.65em}
\textbf{Consulta 5 (1 punto).} Calcula el stock medio por categoría. Redondea la media a dos decimales y ordena por ella de menor a mayor y por categoría.\par\vspace{0.65em}
\textbf{Consulta 6 (1 punto).} Cuenta los pedidos distintos por día del mes y estado, sin distinguir mes ni año. Muestra día, estado y número de pedidos; ordena por día y por número de pedidos descendente.\par\vspace{0.65em}
\textbf{Consulta 7 (2 puntos).} Pista: puedes usar DATEDIFF(). Relaciona cada pedido con su cliente y calcula, por país, el promedio de días entre el registro del cliente y el pedido. Muestra también el número de pedidos y ordena por número de pedidos descendente y país. Redondea el promedio a un decimal.\par\vspace{0.65em}
\textbf{Consulta 8 (2 puntos).} Pista: puedes usar LEAST(). Por año de pedido, muestra el número de pedidos, la suma de todos sus costes y la suma de los costes limitados a 1000 euros por pedido. Ordena por año.\par\vspace{0.65em}
\textbf{Extra (0.5 puntos)} ¿Ha habido algún pago en 29 de febrero? ¿Y algún pedido?
\end{document}
